\documentclass[12pt]{article}
\usepackage[margin=1in]{geometry}
\usepackage{amsmath, amssymb, amsthm}
\usepackage{booktabs}
\usepackage{hyperref}
\usepackage{listings}
\usepackage{xcolor}
\usepackage{graphicx}
\usepackage{enumitem}
\usepackage{fancyhdr}
\usepackage{titlesec}

\hypersetup{colorlinks=true, linkcolor=blue, urlcolor=blue, citecolor=blue}

\lstset{
  language=Python,
  basicstyle=\ttfamily\small,
  keywordstyle=\color{blue},
  commentstyle=\color{gray},
  stringstyle=\color{orange!80!black},
  showstringspaces=false,
  breaklines=true,
  frame=single,
  backgroundcolor=\color{gray!8},
  numbers=left,
  numberstyle=\tiny\color{gray},
  numbersep=6pt
}

\pagestyle{fancy}
\fancyhf{}
\rhead{EE645 -- Physical Models in Remote Sensing}
\lhead{DOM Documentation}
\cfoot{\thepage}

\title{
  \textbf{Discrete Ordinates Method for Leaf-Canopy Radiative Transfer}\\[0.5em]
  \large A Self-Study Guide: Quadrature, Diamond Differences,\\
  Source Iteration, and Extensible Design
}
\author{EE645 -- Physical Models in Remote Sensing}
\date{Spring 2026}

\begin{document}
\maketitle
\tableofcontents
\newpage

% ===========================================================================
\section{The Radiative Transfer Equation}
% ===========================================================================

The radiative transfer equation (RTE) describes how the specific intensity
$I(\mathbf{r}, \boldsymbol{\Omega})$ [W m$^{-2}$ sr$^{-1}$] changes along a
ray traveling in direction $\boldsymbol{\Omega}$:

\begin{equation}
  \frac{dI}{ds} = -\kappa\,I + j
\end{equation}

where $s$ is path length, $\kappa$ is the extinction coefficient, and $j$ is
the emission/scattering source. For a \emph{plane-parallel leaf canopy}
parameterized by the cumulative leaf area index $L$ (0 at canopy top,
$L_\mathrm{AI}$ at ground), the equation becomes

\begin{equation}
  \label{eq:rte}
  \mu\,\frac{\partial I(L,\boldsymbol{\Omega})}{\partial L}
  = -G\,I(L,\boldsymbol{\Omega})
    + \frac{G\,\omega_L}{\pi}\int_{4\pi}
      \Gamma(\boldsymbol{\Omega}'\!\to\!\boldsymbol{\Omega})\,
      I(L,\boldsymbol{\Omega}')\,d\boldsymbol{\Omega}'
\end{equation}

where:
\begin{itemize}[nosep]
  \item $\mu = \cos\theta$ is the polar cosine of direction $\boldsymbol{\Omega}$
        ($\mu < 0$ downward, $\mu > 0$ upward)
  \item $G = 0.5$ for a uniform (spherical) leaf-angle distribution
  \item $\omega_L = \rho_L + \tau_L$ is the leaf single-scattering albedo
  \item $\Gamma(\boldsymbol{\Omega}'\!\to\!\boldsymbol{\Omega})$ is the
        bi-Lambertian scattering phase function [sr$^{-1}$]
\end{itemize}

The integral over all directions couples every direction to every other
direction. This is the fundamental difficulty -- it cannot be solved
analytically in the general case.

\subsection{Bi-Lambertian Phase Function}

The scattering angle $\beta$ between incident direction $(\mu_i, \phi_i)$ and
scattered direction $(\mu_j, \phi_j)$ satisfies:

\begin{equation}
  \cos\beta = \mu_i\mu_j
  + \sqrt{1-\mu_i^2}\,\sqrt{1-\mu_j^2}\cos(\phi_i - \phi_j)
\end{equation}

The bi-Lambertian phase function (Myneni et al.) is:

\begin{equation}
  \Gamma(\boldsymbol{\Omega}_i \to \boldsymbol{\Omega}_j)
  = \frac{\omega_L}{3\pi}\bigl(\sin\beta - \beta\cos\beta\bigr)
    + \frac{\tau_L}{3}\cos\beta
\end{equation}

The first term is diffuse (Lambertian) reflection; the second is forward
transmission through leaves.

% ===========================================================================
\section{Gauss--Legendre Quadrature}
% ===========================================================================

\subsection{The Core Idea}

We replace the continuous directional integral in~\eqref{eq:rte} with a
finite weighted sum:

\begin{equation}
  \int_{-1}^{1} f(\mu)\,d\mu \;\approx\; \sum_{i=1}^{N} w_i\,f(\mu_i)
\end{equation}

The question is: \emph{where do we place the $N$ nodes $\mu_i$, and what
weights $w_i$ do we assign?}

\subsection{Why Legendre Polynomial Zeros?}

\textbf{Fundamental theorem of Gaussian quadrature:} Among all possible
choices of $N$ nodes and $N$ weights, the choice that maximizes the
polynomial degree integrated exactly is achieved if and only if the nodes
$\mu_1, \ldots, \mu_N$ are the zeros of the $N$-th Legendre polynomial
$P_N(\mu)$.

With this choice the rule is \emph{exact} for all polynomials of degree
$\leq 2N - 1$, using only $N$ function evaluations. Uniform spacing can only
achieve exactness up to degree $N - 1$ (the trapezoidal rule).

\subsection{The Legendre Polynomials}

The Legendre polynomials $P_n(\mu)$ are orthogonal on $[-1, 1]$:

\begin{equation}
  \int_{-1}^{1} P_m(\mu)\,P_n(\mu)\,d\mu
  = \frac{2}{2n+1}\,\delta_{mn}
\end{equation}

They are generated by the three-term recurrence:

\begin{align}
  P_0(\mu) &= 1 \\
  P_1(\mu) &= \mu \\
  (n+1)\,P_{n+1}(\mu) &= (2n+1)\,\mu\,P_n(\mu) - n\,P_{n-1}(\mu)
\end{align}

The first few roots (for $N = 4$) are approximately
$\mu \in \{-0.861, -0.340, +0.340, +0.861\}$,
which are symmetric about zero and cluster toward the endpoints -- denser
sampling near horizontal and near-zenith directions.

\subsection{Computing the Weights}

Given a node $\mu_i$ (a zero of $P_N$), the corresponding weight is:

\begin{equation}
  w_i = \frac{2}{(1 - \mu_i^2)\,[P_N'(\mu_i)]^2}
\end{equation}

NumPy computes these via the \emph{Golub--Welsch algorithm}: it constructs the
symmetric tridiagonal matrix whose eigenvalues are the Gauss nodes and whose
eigenvectors give the weights. This is numerically stable for large $N$.

\begin{lstlisting}
# numpy's built-in: returns nodes in ascending order on (-1, 1)
nodes, weights = np.polynomial.legendre.leggauss(N)
\end{lstlisting}

\subsection{Azimuthal Quadrature}

For the azimuthal integral over $\phi \in [0, 2\pi)$, we use $M$ \emph{uniform}
points because the integrand is periodic and smooth:

\begin{equation}
  \phi_m = \frac{(m-1)\,2\pi}{M}, \quad
  w_\phi = \frac{2\pi}{M}, \quad m = 1,\ldots,M
\end{equation}

This is the rectangle rule, which converges exponentially for smooth periodic
functions.

\subsection{Index Convention (1-Based)}

The code uses 1-based arrays with a dummy index-0 slot so that array index
$k$ maps directly to the notation $k = 1, \ldots, N$:

\begin{lstlisting}
mu   = np.zeros(N + 1)   # index 0 unused
mu[1:N+1] = nodes        # nodes in ascending order
# mu[1..N/2]   < 0  (downward)
# mu[N/2+1..N] > 0  (upward)
\end{lstlisting}

% ===========================================================================
\section{The Uncollided / Collided Decomposition}
% ===========================================================================

\subsection{Motivation}

The direct solar beam is a \emph{delta function} in angle -- infinitely
concentrated in one direction. Representing it on a finite quadrature grid
would require enormous $N$ for accuracy. Instead we isolate it analytically:

\begin{equation}
  I(L, \boldsymbol{\Omega}) = I_0(L, \boldsymbol{\Omega})
                             + I_c(L, \boldsymbol{\Omega})
\end{equation}

\begin{description}[nosep]
  \item[$I_0$:] \textbf{Uncollided field.} Radiation that has never scattered.
    No source term -- just exponential attenuation. Solved analytically.
  \item[$I_c$:] \textbf{Collided field.} Radiation that has scattered at least
    once. Smooth in angle -- well-represented on the quadrature grid.
    Driven by the \emph{First Collision Source} $Q$.
\end{description}

\subsection{Uncollided Field: Beer--Lambert}

For the direct solar beam (intensity $I_\mathrm{beam} = f_\mathrm{dir} F_\mathrm{in} / |\mu_\odot|$):

\begin{equation}
  I_0^\mathrm{dir}(L) = I_\mathrm{beam}\,
  \exp\!\left(-\frac{G\,L}{|\mu_\odot|}\right)
\end{equation}

For isotropic diffuse sky (intensity $I_\mathrm{sky} = (1-f_\mathrm{dir})F_\mathrm{in}/\pi$)
along each downward Gauss direction $(\mu_n, \phi_m)$:

\begin{equation}
  I_0(L, \mu_n, \phi_m) = I_\mathrm{sky}\,
  \exp\!\left(-\frac{G\,L}{|\mu_n|}\right), \quad n = 1,\ldots,N/2
\end{equation}

The upward uncollided field comes from Lambertian ground reflection of the
total downward flux reaching the ground:

\begin{equation}
  I_0^\mathrm{gnd} = \frac{\rho_g}{\pi}\,F_\downarrow(L_\mathrm{AI})
\end{equation}

and then attenuates upward through the canopy for each upward Gauss direction.

\subsection{First Collision Source}

When an uncollided photon scatters for the first time it feeds the collided
field. The source has three contributions:

\begin{align}
  Q_1(L, \boldsymbol{\Omega}) &= \frac{1}{\pi}\,
    \Gamma(\boldsymbol{\Omega}_\odot\!\to\!\boldsymbol{\Omega})\,
    I_0^\mathrm{dir}(L)
    \quad\text{(direct solar)} \\[4pt]
  Q_2(L, \boldsymbol{\Omega}) &= \frac{1}{\pi}
    \sum_{n=1}^{N/2}\sum_{m=1}^{M}
    w_n\,w_\phi\,\Gamma(\boldsymbol{\Omega}_{nm}\!\to\!\boldsymbol{\Omega})\,
    I_0(L,\boldsymbol{\Omega}_{nm})
    \quad\text{(diffuse sky)} \\[4pt]
  Q_3(L, \boldsymbol{\Omega}) &= \frac{1}{\pi}
    \sum_{n=N/2+1}^{N}\sum_{m=1}^{M}
    w_n\,w_\phi\,\Gamma(\boldsymbol{\Omega}_{nm}\!\to\!\boldsymbol{\Omega})\,
    I_0(L,\boldsymbol{\Omega}_{nm})
    \quad\text{(ground-reflected)} \\[4pt]
  Q &= Q_1 + Q_2 + Q_3
\end{align}

% ===========================================================================
\section{Diamond Difference: Discretizing the Sweep}
% ===========================================================================

\subsection{Setting Up the 1-D ODE}

After fixing a discrete direction $(\mu_n, \phi_m)$, the RTE reduces to a
1-D first-order ODE in $L$. We divide the canopy into $K$ uniform layers of
thickness $\Delta L = L_\mathrm{AI}/K$. The unknown is the intensity at each
of the $K+1$ layer interfaces (edges):

\begin{equation}
  \mu_n\,\frac{I_{k+1} - I_k}{\Delta L}
  = -G\,\underbrace{\frac{I_k + I_{k+1}}{2}}_{\text{cell-centre}}
    + J_{k+1/2}
\end{equation}

The left side is a finite-difference derivative. The right side evaluates the
source $J$ at the cell centre. The \textbf{Diamond Difference} approximation
is the choice:

\begin{equation}
  \boxed{I_{k+1/2} = \frac{I_k + I_{k+1}}{2}}
\end{equation}

This is the ``diamond'' -- on a mesh plot of space vs.\ angle, the four points
$(I_k, I_{k+1}, I_{k+1/2}, J_{k+1/2})$ form a diamond shape.

\subsection{Deriving the Sweep Equations}

Define the dimensionless optical half-thickness:
\begin{equation}
  f_n = \frac{G\,\Delta L}{\mu_n} \quad\text{(signed: negative for downward)}
\end{equation}

\paragraph{Downward directions ($\mu_n < 0$, sweep from $k=1$ to $k=K$):}

Rearranging the DD equation to solve for $I_{k+1}$ given $I_k$:

\begin{equation}
  \boxed{
  I_{k+1} = a_n\,I_k - b_n\,J_{k+1/2}
  }
  \qquad
  a_n = \frac{1+(1-\alpha)f_n}{1-\alpha f_n},
  \quad
  b_n = \frac{f_n}{G\,(1-\alpha f_n)}
\end{equation}

Since $\mu_n < 0 \Rightarrow f_n < 0$, we get $a_n < 1$ (attenuation) and
$-b_n > 0$ (source adds to intensity). Both are physically correct.

\paragraph{Upward directions ($\mu_n > 0$, sweep from $k=K$ to $k=1$):}

\begin{equation}
  \boxed{
  I_k = c_n\,I_{k+1} + d_n\,J_{k+1/2}
  }
  \qquad
  c_n = \frac{1-(1-\alpha)f_n}{1+\alpha f_n},
  \quad
  d_n = \frac{f_n}{G\,(1+\alpha f_n)}
\end{equation}

Since $\mu_n > 0 \Rightarrow f_n > 0$, we get $c_n < 1$ and $d_n > 0$.

\subsection{The Role of $\alpha$}

The parameter $\alpha \in [0, 1]$ controls how the cell-centre value is
weighted between the two edges:

\begin{equation}
  I_{k+1/2} = (1-\alpha)\,I_k + \alpha\,I_{k+1}
\end{equation}

\begin{center}
\begin{tabular}{cll}
  \toprule
  $\alpha$ & Name & Properties \\
  \midrule
  $0.5$ & Diamond Difference & 2nd-order accurate; may produce negative $I$ \\
  $1.0$ & Step (upwind) & 1st-order accurate; always positive \\
  \bottomrule
\end{tabular}
\end{center}

\subsection{Sign Convention: Why Signed $\mu$ Matters}

Using the \emph{signed} cosine in $f_n = G\Delta L/\mu_n$ (not $|\mu_n|$) is
critical. For a downward direction $\mu_n = -0.5$:

\begin{align*}
  f_n &= G\,\Delta L / (-0.5) < 0 \\
  a_n &= \frac{1 - 0.5|f_n|}{1 + 0.5|f_n|} < 1 \quad\checkmark \text{ (attenuating)}\\
  \text{source term} &= -b_n\,J > 0 \quad\checkmark \text{ (source adds intensity)}
\end{align*}

If instead $|\mu_n|$ were used, $f_n > 0$ for downward, giving $a_n > 1$
(amplifying -- unphysical) and the source term would \emph{subtract} from $I$
(wrong sign). This is a common implementation bug.

\subsection{The Fix-Up}

Standard DD can produce negative intensities in optically thick cells (large
$G\,\Delta L / |\mu_n|$). Negative intensity is unphysical, so we clip:

\begin{lstlisting}
I[k+1] = max(a * I[k] - b * J[k], 0.0)
\end{lstlisting}

The proper solution is to use more layers (larger $K$) so that each cell is
optically thin. As $\Delta L \to 0$, DD becomes exact.

% ===========================================================================
\section{Source Iteration}
% ===========================================================================

The collided field equation contains a self-scattering term -- $I_c$ depends
on itself through the multiple-scattering source $S$:

\begin{equation}
  S(L,\boldsymbol{\Omega}) = \frac{1}{\pi}
  \sum_{n,m} w_n\,w_\phi\,\Gamma(\boldsymbol{\Omega}_{nm}\!\to\!\boldsymbol{\Omega})\,
  I_c(L,\boldsymbol{\Omega}_{nm})
\end{equation}

We solve this by \textbf{Source Iteration} (also called $\Lambda$-iteration):

\begin{enumerate}
  \item Initialize: $I_c = 0$, $S = 0$
  \item Compute total source: $J = Q + S$
  \item Sweep all $N \times M$ directions using DD to obtain new $I_c$
  \item Compute new $S$ from $I_c$
  \item Check convergence:
    \begin{equation}
      \epsilon = \max_k
      \frac{|SI_k^\mathrm{new} - SI_k^\mathrm{old}|}
           {\max(|SI_k^\mathrm{new}|,\;\varepsilon_0)}
      < \mathrm{TOL}
    \end{equation}
    where $SI_k = \sum_{n,m} w_n\,w_\phi\,I_{c,k+1/2}^{n,m}$ is the
    scalar irradiance (integral of intensity over all directions, no cosine weight)
  \item Repeat from step 2
\end{enumerate}

\subsection{Convergence Rate}

Source iteration converges because each scatter removes energy ($\omega_L < 1$).
The spectral radius of the iteration is approximately $\omega_L$:

\begin{center}
\begin{tabular}{ccc}
  \toprule
  Band & $\omega_L$ & Typical iterations \\
  \midrule
  RED ($\rho_L=0.06, \tau_L=0.04$) & 0.10 & 3--5 \\
  NIR ($\rho_L=0.525, \tau_L=0.450$) & 0.975 & 100--500 \\
  \bottomrule
\end{tabular}
\end{center}

For NIR, the slow convergence is well-known and can be accelerated with
\emph{Diffusion Synthetic Acceleration (DSA)} -- beyond the scope of this course.

\subsection{Ground Boundary Condition (Collided)}

The downward sweep is completed first. The resulting collided downward flux at
the ground is immediately reflected Lambertianly to set the lower boundary
condition for the upward sweep \emph{in the same iteration}:

\begin{equation}
  I_c^\mathrm{gnd} = \frac{\rho_g}{\pi}
  \sum_{n=1}^{N/2}\sum_{m=1}^{M}
  w_n\,w_\phi\,|\mu_n|\,I_c(\mu_n,\phi_m,L_\mathrm{AI})
\end{equation}

Using the current iteration's downward IC (rather than the previous iteration's)
for the ground BC significantly accelerates convergence.

% ===========================================================================
\section{Energy Balance}
% ===========================================================================

At convergence, the total irradiances are:

\begin{align}
  F_\downarrow^\mathrm{tot}(L) &= F_\downarrow^\mathrm{dir}(L)
    + F_\downarrow^\mathrm{dif}(L) + F_\downarrow^\mathrm{col}(L) \\
  F_\uparrow^\mathrm{tot}(L) &= F_\uparrow^\mathrm{unc}(L)
    + F_\uparrow^\mathrm{col}(L)
\end{align}

The correct energy balance is:

\begin{equation}
  \boxed{
  F_\mathrm{in} = F_\uparrow^\mathrm{tot}(0)
    + A_\mathrm{leaves}
    + (1 - \rho_g)\,F_\downarrow^\mathrm{tot}(L_\mathrm{AI})
  }
\end{equation}

where the leaf absorption is:

\begin{equation}
  A_\mathrm{leaves} = G\,(1-\omega_L)\,\Delta L
  \sum_{k=1}^{K} SI_{k+1/2}
\end{equation}

and $SI_{k+1/2}$ is the scalar irradiance (intensity integrated over all
directions without a cosine weight):

\begin{equation}
  SI_{k+1/2} = I_0^\mathrm{dir}{}_{k+1/2}
  + \sum_{n,m} w_n\,w_\phi
    \bigl[I_{0,k+1/2}^{n,m} + I_{c,k+1/2}^{n,m}\bigr]
\end{equation}

\textbf{Note:} The balance uses $(1-\rho_g)\,F_\downarrow^\mathrm{tot}$ at the
ground (ground-absorbed flux), not the full $F_\downarrow^\mathrm{tot}$. The
ground-reflected portion $\rho_g\,F_\downarrow^\mathrm{tot}$ is already
accounted for in $F_\uparrow^\mathrm{tot}$ and $A_\mathrm{leaves}$.

% ===========================================================================
\section{Alternative Methods and Extensible Design}
% ===========================================================================

\subsection{Alternative Quadrature Schemes}

\begin{description}
  \item[Double Gauss (recommended upgrade):] Apply separate Gauss--Legendre
    quadrature on each hemisphere $[-1,0]$ and $[0,1]$ independently, then
    combine. Each half-hemisphere integral of degree $\leq 2(N/2)-1$ is exact.
    This gives significantly better hemisphere flux accuracy for the same $N$,
    and is used in production codes such as DISORT.

    \emph{Implementation:} Apply \texttt{leggauss(N/2)} on $[0,1]$ to get
    upward nodes, then mirror to get downward nodes.

  \item[Gauss--Lobatto:] Nodes include the endpoints $\pm 1$ (nadir and
    zenith). Useful when boundary behaviour dominates, but sacrifices one
    order of accuracy.

  \item[Uniform (equal-weight):] Evenly spaced $\mu_i = -1 + (2i-1)/N$ with
    equal weights $w_i = 2/N$. Simple but only 1st-order accurate. Use only
    for debugging.
\end{description}

\subsection{Alternative Sweep Schemes}

\begin{center}
\begin{tabular}{llll}
  \toprule
  Scheme & Accuracy & Always positive? & Notes \\
  \midrule
  Diamond Difference ($\alpha=0.5$) & 2nd order & No & Standard; needs fix-up \\
  Step / Upwind ($\alpha=1.0$) & 1st order & Yes & Same code, change $\alpha$ \\
  Linear Discontinuous (LD) & 2nd order & Yes & 2$\times$ unknowns per cell \\
  \bottomrule
\end{tabular}
\end{center}

\paragraph{Linear Discontinuous Galerkin (LD):} Each cell has \emph{two}
unknowns -- a left-edge and a right-edge value -- related by a linear
interpolation. The scheme is always positive and 2nd-order accurate. The
tradeoff is that the inner loop must solve a $2\times2$ system per cell
rather than a scalar recurrence.

\subsection{Suggested Modular Design}

When rewriting from scratch, separate concerns into distinct modules:

\begin{lstlisting}[language=bash,numbers=none]
DOM_v2/
  config.py       <- all inputs + method flags (QUADRATURE, SWEEP, ALPHA_DD)
  quadrature.py   <- get_polar_quadrature(N, method)
  phase.py        <- gamma_scalar(), precompute_G_qq(), precompute_G_sol()
  uncollided.py   <- solve_uncollided(), compute_Q()
  collided.py     <- solve_collided() with pluggable sweep_coeffs()
  energy.py       <- scalar_irradiance(), energy_balance()
  brf.py          <- brf_at_view() via DOM interpolation
  main.py         <- orchestrates everything, prints diagnostics
\end{lstlisting}

Key design principle: \texttt{config.py} holds \emph{all} parameters. Every
function receives only what it needs. This makes it easy to swap quadrature
method or sweep scheme without touching solver logic.

% ===========================================================================
\section{Quick Reference: Key Equations}
% ===========================================================================

\begin{center}
\renewcommand{\arraystretch}{1.6}
\begin{tabular}{ll}
  \toprule
  Quantity & Formula \\
  \midrule
  Legendre zeros (nodes) & zeros of $P_N(\mu)$; computed via \texttt{leggauss(N)} \\
  Gauss weights & $w_i = 2\,/\,[(1-\mu_i^2)(P_N'(\mu_i))^2]$ \\
  Azimuthal weight & $w_\phi = 2\pi/M$ \\
  Uncollided (direct) & $I_0^\mathrm{dir}(L) = I_\mathrm{beam}\,e^{-GL/|\mu_\odot|}$ \\
  Uncollided (diffuse) & $I_0(L,\mu_n) = I_\mathrm{sky}\,e^{-GL/|\mu_n|}$ \\
  DD coefficients & $f = G\Delta L/\mu$;\; $a=(1+(1{-}\alpha)f)/(1-\alpha f)$;\;
                    $b=f/[G(1-\alpha f)]$ \\
  DD sweep (down) & $I_{k+1} = a\,I_k - b\,J_{k+1/2}$ \\
  DD sweep (up) & $I_k = c\,I_{k+1} + d\,J_{k+1/2}$ \\
  Multiple-scatter source & $S_{ij} = \pi^{-1}\sum_{nm}w_n w_\phi\,\Gamma_{nm\to ij}\,I_{c,nm}$ \\
  Scalar irradiance & $SI_k = I_0^\mathrm{dir}{}_k + \sum_{nm}w_n w_\phi(I_{0,nm,k}+I_{c,nm,k})$ \\
  Leaf absorption & $A = G(1-\omega_L)\Delta L\sum_k SI_k$ \\
  Energy balance & $F_\mathrm{in} = F_\uparrow(0) + A + (1-\rho_g)F_\downarrow(L_\mathrm{AI})$ \\
  \bottomrule
\end{tabular}
\end{center}

\end{document}
